\chapter{Surfaces}

Surfaces describe loadable or connectable element sides. They are used for pressure loads, distributed surface loads, tie constraints, surface couplings, and node-to-surface contact. A surface is a selection of element sides; it is not a material or section property.

\section{Surface from one element side}

\begin{syntax}
*SURFACE, NAME=<surface-set>
<surface-id>, <element-id>, <side>
\end{syntax}

The side can be written as a number or as a token such as \val{S1} or \val{S2}. For shells, \val{SPOS} and \val{SNEG} select the positive or negative shell side.

\section{Surface from an element set}

\begin{syntax}
*SURFACE, NAME=<surface-set>, TYPE=ELEMENT
<element-set-or-id>, <side>
\end{syntax}

If the target is an element set, FEMaster creates surfaces for the selected side of all elements in the set. Otherwise the target is interpreted as a single element ID.

\section{Orientation and sign}

Surface loads and contact depend on the chosen side and its normal direction. The normal follows the local face orientation implied by the element connectivity and selected side. If pressure or traction direction is unexpected, check the selected side, element orientation, and shell side. For shells, positive and negative sides share the same geometric midsurface but have opposite normals.

Surface orientation is particularly important for \cmd{CONTACT}. The node-to-surface formulation evaluates the signed gap with the current master-face normal at the bounded closest point. A master surface that is oriented opposite to the intended contact direction can be corrected for the contact pair with \key{FLIP=YES}; this does not change the underlying element connectivity or surface definition.

\begin{warningbox}[Contact surfaces]
The current \cmd{CONTACT} interface uses a master and a slave surface set. The unique nodes of the slave surface set become the contact nodes, while the master surface set provides the candidate master faces. There is no mortar, surface-to-surface integration, or user-facing search-distance parameter.
\end{warningbox}
